\section{本章の目的と位置づけ}
\label{sec:sim_purpose}

前章では，軟弱地面における床沈み込みが，重心（CoM）高さ変動や ZMP 応答を通じて，従来の歩行制御手法に影響を与える可能性があることを，力学的観点から整理した．しかし，これらの議論は理論的考察に基づくものであり，実際の歩行中にどのような挙動が生じているかについては，必ずしも明らかではない．

そこで本章では，軟弱地面を模擬したシミュレーション環境を構築し，従来手法を用いた歩行において実際に観測される挙動を整理する．本章の目的は，軟弱地面において歩行が不安定となる現象を定量的に観測し，その特徴を明らかにすることである．

なお，本章では不安定化の原因やその解決手法については扱わず，あくまでシミュレーション結果に基づく挙動の整理に焦点を当てる．
